\section{Fazit und Ausblick}

Die Projektarbeit zeigt für den kontrollierten Alurohrdatensatz die technische
Machbarkeit einer Docker-basierten Scan-to-BIM-Pipeline. SAM~3.1, unmaskiertes
COLMAP-SfM, ideale Bilddomäne, objektspezifisches STS, Original-GS-
Meshgewinnung und DGtal-Centerline können im normalen Inline-Vollauf und im
Matrixbetrieb als reproduzierbare Kette betrieben und archiviert werden. Der
Inline-Vollauf führt den Maskenwarp und die zugehörige Coverage-Prüfung
ebenfalls vor STS aus. Dies ist durch die archivierten Autopilot-Vollläufe
nachgewiesen (Abschnitt~\ref{sec:versuchsaufbau}).

Der Eigenbeitrag besteht insbesondere in der kameramodellbewussten Trennung von
Roh- und Idealbilddomäne, der konsistenten Nutzung der SAM-Masken, den
prüfbaren Datenverträgen zwischen fünf Containerumgebungen, der direkten
Original-GS-Meshroute sowie der automatisierten Versuchs- und
Fehlerauswertung. Interaktiver CLI-Betrieb, Autopilot und Matrixrunner decken
Erstlauf, Standardlauf und Batchverarbeitung ab.

Variante A ist für den Projektstand die bevorzugte Meshroute, weil sie die
STS-Gaussian-Geometrie direkter erhält als eine zusätzliche SuGaR-Coarse-
Optimierung. PINHOLE und OPENCV sind als Kameramodell-Ablationen technisch
vergleichbar. SuGaR-Coarse bleibt eine wichtige Vergleichsroute.

Die Ergebnisse liefern einen Funktionsnachweis, aber keinen
Vermessungsnachweis. Die Bachelorarbeit muss deshalb reale Aufnahmen,
geeignete GCP-/GNSS-Messungen, echte 4x4-Georeferenzierung und geometrische
Fehlermetriken ergänzen. Die in dieser Arbeit aufgebaute Struktur trennt diese
wissenschaftlichen Fragen von den bereits geprüften technischen
Zwischenschritten. Insbesondere sind die von COLMAP geschätzten
Kameraverzeichnungsparameter ohne unabhängige Checkerboard- oder
ChArUco-Kalibrierung nicht als physikalisch validierte Objektivdaten zu
interpretieren.
